% Part: lambda-calculus
% Chapter: introduction
% Section: basic-pr-lambda

\documentclass[../../../include/open-logic-section]{subfiles}

\begin{document}

\olfileid{lam}{int}{bas}
\olsection{الدوال العودية البدائية الأساسية \usetoken{S}{lambda definable}}

\begin{lem}
كل واحدة من الدوال~$\Zero$ و$\Succ$ و$\Proj{n}{i}$ هي !!{lambda definable}.
\end{lem}

\begin{proof}
أما~$\Zero$ فيمثلها الحد
$\lambd[x][\lambd[y][y]]$.

وأما دالة الخلف~$\Succ$ فتعريفها
$\fn{Succ}(u) = \lambd[x][\lambd[y][x(uxy)]]$. ولتتأمل وجه هذا التعريف:
لكل عدد~$\num{n}$ نعدّه مكرّرًا، ولكل دالة~$f$، تكون
$\fn{Succ}(\num{n},f)$ دالة تأخذ~$y$، فتطبق~$f$
$n$ مرة، بادئةً من~$y$، ثم تزيد تطبيقًا واحدًا.

وأما الإسقاطات فأمرها بيّن:
$\fn{Proj}^n_i(x_0, \dots, x_{n-1}) = x_i$. فمقتضى اصطلاحنا أن
$\fn{Proj}^n_i$ هي حد لامبدا
$\lambd[x_0][\dots \lambd[x_{n-1}][x_i]]$.
\end{proof}

\end{document}

